\subsection{Factoriality of rings of formal power series}

PROPOSITION 8. Let C be a ring which is either a field or a discrete valuation ring. Then the domain of formal power series C{[}{[}X₁, . . ., Xₙ{]}{]} is factorial.

Let p be the maximal ideal of C and \(\pi\) a generator of p (if C is a field, then \(\pi = 0\) ). Let C be given the p-adic topology, which is Hausdorff. As C is a Noetherian local ring, \(B = C[[X_{1}, \ldots, X_{n}]]\) is a Noetherian local ring and its completion is \(\hat{C}[[X_{1}, \ldots, X_{n}]]\) (Chapter III, §2, no. 6, Proposition 6). By the Corollary to Proposition 4 (no. 6), it suffices to prove that \(\hat{C}[[X_{1}, \ldots, X_{n}]]\) is factorial. Now, if C is a field, then \(\hat{C} = C\) ; if C is a discrete valuation ring, the same is true of \(\hat{C}\) (Chapter VI, §5, no. 3, Proposition 5). We shall therefore assume in the remainder of the proof that C is complete.

Arguing by induction starting with the trivial case n = 0, we shall assume that it has been proved that \(A = C[[X_{1}, \ldots, X_{n-1}]]\) is factorial. We shall identify B with \(A[[X_{n}]]\) and denote by m the maximal ideal of A (generated by \(\pi\) , \(X_{1}, \ldots, X_{n-1}\) ). We shall prove that every non-zero element g of B is, in an essentially unique way, a product of extremal elements.

Let K be the field C/Cπ; as B/Bπ is identified with K{[}{[}X₁, . . . , Xₙ{]}{]}, the ideal Bπ is prime and π is extremal. If π ≠ 0, \(B_{B\pi}\) is therefore the ring of a normed discrete valuation w (Chapter VI, § 3, no. 6, Proposition 9); every nonzero element g of B may therefore be written as g = πw(g)f, where f ∈ B and f is not a multiple of π. It will therefore suffice to show that f is an essentially unique product of extremal elements. Now the canonical image of f in K{[}{[}X₁, . . . , Xₙ{]}{]} is not zero; Lemma 3 (no. 7) therefore shows that there exists an automorphism of B which maps f to an element f' such that the coefficients of f'(0, . . . , 0, Xₙ) are not all in Cπ; this means that the coefficients of the series f', considered as a formal power series in Xₙ, are not all in m. It will suffice to prove our assertion for f'.

In what follows, all the elements of B will be considered as formal power series in \(X_{n}\) with coefficients in A. By Proposition 6 of no. 8 (applicable since C and therefore A are separable and complete and the reduced series of \(f'\) is \(\neq0\) ), \(f'\) is associated, in B, with a unique distinguished polynomial F. By Proposition 7 of no. 8, every series which divides \(f'\) (or, what amounts to the same, which divides F) is associated with a distinguished polynomial which divides F and every decomposition of \(f'\) is, to within invertible factors, of the form \(f' = uF_{1}\ldots F_{q}\) , where u is invertible and the \(F_{t}\) are extremal distinguished polynomials (in B) such that \(F = F_{1}\ldots F_{q}\) . By the Corollary to Proposition 7 of no. 8, the \(F_{t}\) are also extremal in \(A[X_{n}]\) . Now, as A is factorial by the induction hypothesis, so is \(A[X_{n}]\) (Theorem 2, no. 5); hence, since they are monic, the \(F_{i}\) are uniquely determined by F (up to a permutation). This shows the uniqueness of the decomposition \(f' = uF_{1} \ldots F_{q}\) ; its existence follows from the fact that B is Noetherian, which completes the proof.

\textbf{Remarks}\par

\begin{enumerate}
\def\labelenumi{(\arabic{enumi})}
\item
  There exist factorial rings A such that the ring A{[}{[}X{]}{]} is not factorial (Exercise 8). However, if A is a principal ideal domain, A{[}{[}X₁, \ldots, Xₙ{]}{]} is factorial (Exercise 9).
\item
  * We shall see later, by homological methods, that every regular local ring is factorial (cf.~§ 4, no. 7, Corollary 3 to Proposition 16). This will give another proof, conceptually simpler, of Proposition 8. *
\end{enumerate}

